Наименование работы: Разработка веб-приложения для администрирвания групповых занятий

Цель работы: Создать интуитивно понятную платформу, которая объединит функции планирования, учета,
финансового контроля и коммуникации, адаптированную под потребности
преподавателей и частных образовательных учреждений

Задачи для достижения цели:

\begin{enumerate}
    \item Проанализировать предметную область;
    \item Проанализировать целевую аудиторию;
    \item Проанализировать конкурентов и аналоги;
    \item Определить требования к системе;
    \item Спроектировать и описать бизнес процессы;
    \item Спроектировать и разработать веб-приложение;
    \item Спроектировать пользовательский интерфейс;
    \item Протестировать систему;
\end{enumerate}
Объект и предмет исследования: Веб-приложение

Работа состоит из 3 глав, Введения, Заключения и Приложений. Общий объем работы …. страниц, из них … приложений на … листах. В работу включены … рисунков, …. таблиц. Библиографический список состоит из …. источников, из них …. печатных источников, …. интернет-источников.

Во Введении описаны цель и задачи работы, объект и предмет исследования, его актуальность, новизна и практическая значимость.

В Первой главе описывается аналитическая часть разработки. Исследована предметная область, целевая аудитория и требования к системе

Во Второй главе представлена разработка серверной и клиентской части веб-приложения.

В Третьей главе раскрыты технические аспекты внедрения и опытной эксплуатации полученного практического результата ВКР.

В Заключении представлены выводы по поставленным задачам.